% =====================================================================
%  Apresentação (Beamer) — Projeto Hermathena / BINGO
%  Compilar com LuaLaTeX ou XeLaTeX (Overleaf: Menu > Compiler).
%  Mantenha hermathena.sty no mesmo diretório.
% =====================================================================
\documentclass[aspectratio=169,11pt]{beamer}

\usepackage{hermathena}
\usetheme{default}
\usefonttheme{professionalfonts}
\setbeamertemplate{navigation symbols}{}

% --- Paleta aplicada ao Beamer --------------------------------------
\setbeamercolor{normal text}{fg=hermaGrafite}
\setbeamercolor{frametitle}{fg=hermaGrafite,bg=hermaBranco}
\setbeamercolor{title}{fg=hermaGrafite}
\setbeamercolor{subtitle}{fg=hermaAzulBingo}
\setbeamercolor{structure}{fg=hermaAzulBingo}
\setbeamercolor{itemize item}{fg=hermaAzulBingo}
\setbeamercolor{itemize subitem}{fg=hermaLaranja}
\setbeamercolor{block title}{fg=white,bg=hermaAzulProfundo}
\setbeamercolor{block body}{fg=hermaGrafite,bg=hermaBranco}

% --- Fundo e filete ------------------------------------------------
\setbeamercolor{background canvas}{bg=white}
\setbeamertemplate{frametitle}{%
  \vspace{6pt}%
  \usebeamerfont{frametitle}\insertframetitle\par%
  \vspace{-2pt}{\color{hermaLaranja}\rule{\linewidth}{1pt}}%
}

% --- Rodapé --------------------------------------------------------
\setbeamertemplate{footline}{%
  \hspace{1em}{\scriptsize\color{hermaGrafite} Projeto Hermathena · BINGO}%
  \hfill{\scriptsize\color{hermaGrafite}\insertframenumber}\hspace{1em}\vspace{4pt}%
}

\title{Projeto Hermathena}
\subtitle{Comunicação Multimídia do Radiotelescópio BINGO}
\author{Luciano Barosi}
\date{9 de outubro de 2026}

\begin{document}

% --- Capa ----------------------------------------------------------
\begin{frame}[plain]
  \vfill
  \begin{center}
    \hermathenaLogo[3cm]\par\vspace{10pt}
    {\hermaDisplay\color{hermaGrafite}\Huge\bfseries \inserttitle}\par\vspace{4pt}
    {\hermaDisplay\color{hermaAzulBingo}\large \insertsubtitle}\par\vspace{14pt}
    \hermaRule[0.6\linewidth]\par\vspace{8pt}
    {\color{hermaGrafite}\insertauthor\ · \insertdate}
  \end{center}
  \vfill
\end{frame}

% --- Seção ---------------------------------------------------------
\section{Fundamentos}
\begin{frame}{Uma jornada de transformação humana nas asas da ciência}
  \begin{itemize}
    \item Aproximar ciência e sociedade.
    \item Transformar conhecimento em oportunidade.
    \item Interiorizar a ciência no sertão paraibano.
  \end{itemize}
  \vspace{8pt}
  \begin{block}{Princípio}
    Comunicar ciência de fronteira com sóbria formalidade e rigor científico.
  \end{block}
\end{frame}

\begin{frame}{Núcleos do projeto}
  \begin{columns}[T]
    \begin{column}{0.5\textwidth}
      \textbf{Magellan} — infraestrutura digital\\
      \textbf{Triangulum} — produção de conteúdo
    \end{column}
    \begin{column}{0.5\textwidth}
      \textbf{Via Láctea} — educação\\
      \textbf{Andromeda} — engajamento
    \end{column}
  \end{columns}
\end{frame}

\end{document}
